\newpage
\section*{Введение}
\addcontentsline{toc}{section}{Введение}

В современном мире, где объёмы обрабатываемых данных постоянно
растут, а вычислительные задачи становятся всё сложнее, технологии
высокопроизводительных вычислений (HPC) приобретают критически
важное значение. Их применение позволяет объединять ресурсы
множества вычислительных узлов в единый кластер, обеспечивая тем
самым существенное увеличение производительности и сокращение
времени решения сложных научных и инженерных задач.

Эффективная эксплуатация вычислительного кластера невозможна без
системы управления заданиями и ресурсами, обеспечивающей приём
заявок пользователей, их постановку в очередь и распределение по
вычислительным узлам в соответствии с заданной политикой
планирования. Одной из наиболее распространённых систем такого
класса является SLURM (Simple Linux Utility for Resource
Management), широко используемая как в коммерческих, так и в
научно-исследовательских вычислительных центрах.

Освоение принципов администрирования и эксплуатации
вычислительного кластера на базе SLURM не требует обязательного
наличия физического кластера: для учебных и исследовательских
целей вычислительная среда может быть воспроизведена в виде набора
виртуальных машин. Такие инструменты, как VirtualBox, позволяют
создавать изолированные виртуальные среды, в которых каждая
машина выполняет роль отдельного узла кластера, что делает
возможным изучение архитектуры распределённой системы средствами
одного физического компьютера.

Операционные системы семейства Linux, в частности дистрибутив
Fedora Server, благодаря своей гибкости, открытости и широкому
набору системных утилит, предоставляют все необходимые
инструменты для организации управляющего и вычислительных узлов
кластера в режиме интерфейса командной строки. Для организации
безопасного взаимодействия между узлами кластера и удалённого
администрирования незаменим протокол SSH (Secure Shell),
обеспечивающий авторизацию по криптографическим ключам без
передачи пароля по сети, а также служба munge, отвечающая за
аутентификацию сообщений между демонами SLURM.

Помимо развёртывания и настройки самой вычислительной среды,
важной задачей является имитация реальной пользовательской
нагрузки на кластер с целью проверки корректности работы
диспетчера и последующего анализа эффективности планирования
заданий. Для построения такого имитатора целесообразно опираться
не на произвольные, а на статистически обоснованные параметры
заявок, полученные путём анализа реальных данных учёта
выполненных задач (accounting data) действующего вычислительного
кластера. Применение математического аппарата теории вероятностей
и математической статистики, в частности нормального и
логнормального распределений, позволяет генерировать синтетические
заявки, по своим временным характеристикам приближенные к
наблюдаемым в реальных системах.

В ходе настоящей НИР рассматриваются технологии организации
вычислительного кластера на базе виртуализации и Linux, включая
развёртывание управляющего и вычислительных узлов средствами
VirtualBox, настройку сетевого взаимодействия и SSH-доступа,
установку и конфигурирование диспетчера SLURM, а также разработку
имитатора пользовательских заявок, основанного на статистическом
анализе набора учётных данных кластера.